% ===================================================================
%  TABELO N° 2 — La Homala Korpo. — La Amuztempo. La Ludi.
%  Feuillets 13 a 18 du fac-simile = folios 11 a 16.
%  Correspondance : folio imprime = numero de feuillet - 2.
%
%  Les marques « %%K » ne sont pas des commentaires ordinaires : elles
%  portent la CLE D'ALIGNEMENT qui apparie, dans la page de lecture, ce
%  bloc-ci et son homologue francais. La cle se lit
%      t<tableau>-<numero d'alinea numerote>-<rang dans l'alinea>
%  et ne depend d'aucune pagination : les deux editions ne coupent pas
%  leurs pages aux memes endroits, mais elles numerotent les memes
%  alineas — c'est le seul point d'ancrage que les deux livrets
%  partagent, et il est de l'auteur, non du transcripteur.
% ===================================================================

% --- Feuillet 13 (folio 11 : ouverture de tableau) ------------------
\begin{VUpage}[13]{11}
%%K t02-apar-1 apar
\VUsaut{4.0mm}
\VUtitre{15.0pt}{60}{TABELO N\textsuperscript{o} 2}
\VUsaut{2.0mm}
\VUtitre{11.4pt}{0}{La Homala Korpo. --- La Amuztempo.}
\VUtitre{11.4pt}{0}{La Ludi.}
%  « La Korpo homala. » et son sous-titre ne sont pas l'apparat du
%  tableau : c'est le titre de la premiere section, en 13.2pt sous un
%  titre de tableau en 11.4pt, et le volume francais lui repond par
%  « I. Le Corps Humain. (Simple leçon d'histoire naturelle). ». D'ou
%  une cle propre, pour que les deux colonnes se repondent.
%%K t02-tit-0 sub
\VUtitre{13.2pt}{0}{La Korpo homala.}
\VUsaut{2.4mm}
\VUcentre{10.2pt}{0}{\textit{(Simpla leciono pri naturcienco.)}}
\VUsaut{3.0mm}

%  L'imprimeur pose « LA KAPO. » avant l'alinea 1, qui annonce les trois
%  parties du corps et ne traite pas encore de la tete ; la tete
%  commence a l'alinea 2. Le PDF garde la ligne ou elle est imprimee ;
%  « pos= » ne deplace que la page de lecture.
%%K t02-tit-1 sub pos=t02-01-1
\VUpk{10.2pt}{40}{La Kapo.}
\VUsaut{3.0mm}

%%K t02-01-1 p
\VUblancAlinea
1. --- La precipua parti dil korpo homala\nl
esas : la \VUgras{kapo,} la \VUgras{torso} e la \VUgras{membri.}

%%K t02-02-1 p
\VUblancAlinea
2. --- La kapo inkluzas du parti : la \VUgras{kranio,}\nl
qua kontenas la \VUgras{cerebro} e quan kovras la \VUgras{ha}\cc
\VUgras{raro,} e la \VUgras{vizajo.}

%%K t02-03-1 p
\VUblancAlinea
3. --- Sur nia desegnuro ni vidas nek la kra\cc
nio nek la cerebro; ma ni tre klare vidas la im\cc
portanta parti dil vizajo.

%%K t02-04-1 p
\VUblancAlinea
4. --- Yen unesme la \VUgras{fronto,} qua indikas po\cc
tenta intelekto, pensado sempre aganta. La\nl
\VUgras{tempori} esas tre libera. L'\VUgras{okulo} esas vivaca e\nl
larje apertita. Densa \VUgras{brovo} e longa \VUgras{cilii} shir\cc
mas lu.

%%K t02-05-1 p
\VUblancAlinea
5. --- Yen pluse la \VUgras{nazo} e la \VUgras{naztrui.} La nazo\nl
servas ni por flarar e por respirar. Singlalatere\nl
dil nazo esas la \VUgras{vangi,} plu fore l'\VUgras{oreli.} L'in\cc
frajo dil vizajo konsistas ek la \VUgras{boko} e la \VUgras{men}\cc
\VUgras{tono.}

%%K t02-06-1 p
\VUblancAlinea
6. --- Ni ne tre bone vidas li, nam la \VUgras{labio}\cc
\VUgras{barbo} e la \VUgras{vangobarbo} celas li de ni. Ma ni\nl
savas, ke la boko havas la du \VUgras{labii,} la \VUgras{maxilo}\nl
e la \VUgras{mandibulo} kun lia \VUgras{denti,} e la \VUgras{lango.} Per
\parplein
\end{VUpage}

% --- Feuillet 14 (folio 12) ------------------------------------------
\begin{VUpage}[14]{12}
%%K t02-06-1 p suite
\VUcontinue
la boko ni parolas e manjas. Per la lango ni\nl
gustas la nutrivi.

%%K t02-07-1 p
\VUblancAlinea
7. --- L'okulo, l'orelo, la nazo e la boko esas\nl
same kam la fingri, l'organi di la kin sensi : la\nl
\VUgras{vidosenso,} l'\VUgras{audosenso,} la \VUgras{flarosenso,} la \VUgras{gus}\cc
\VUgras{tosenso,} la \VUgras{tushosenso.}

%%K t02-08-1 p
\VUblancAlinea
8. --- La kapo esas do un de la maxim\nl
interesanta parti dil korpo homala.

%%K t02-tit-2 sub
\VUsaut{4.4mm}
\VUpk{10.2pt}{40}{La Torso.}
\VUsaut{3.0mm}

%%K t02-09-1 p
\VUblancAlinea
9. --- La \VUgras{kolo} juntas la kapo a la torso. Ol\nl
inkluzas la \VUgras{guturo} e la \VUgras{nuko.}

%%K t02-10-1 p
\VUblancAlinea
10. --- La torso anke inkluzas multa organi\nl
esencala. En la \VUgras{pektoro} esas la \VUgras{pulmoni,} per\nl
qui ni respiras, e la \VUgras{kordio} qua esas la centro\nl
di la vivo. Kande la kordio cesas pulsar, la\nl
vivanta ento mortas.

%%K t02-11-1 p
\VUblancAlinea
11. --- La \VUgras{kosti} mantenas la pektoro.

%%K t02-12-1 p
\VUblancAlinea
12. --- La cetera parti dil torso esas la\nl
\VUgras{flanko,} la \VUgras{dorso,} la \VUgras{shultri} e l'\VUgras{abdomino} (ven\cc
tro). Ni kushas ni sur la flanko o sur la dorso\nl
por dormar, ni portas la kargaji sur nia shul\cc
tro.

%%K t02-13-1 p
\VUblancAlinea
13. --- En l'abdomino esas la \VUgras{stomako} e la\nl
\VUgras{intestini.} Li servas ni por digestar la nutrivi.

%%K t02-tit-3 sub
\VUsaut{4.4mm}
\VUpk{10.2pt}{40}{La Membri.}
\VUsaut{3.0mm}

%%K t02-14-1 p
\VUblancAlinea
14. --- Ni havas quar \VUgras{membri} : du \VUgras{brakii} e\nl
du \VUgras{gambi.} Per la brakii ni prenas, tenas, levas,\nl
rekolias la objekti; per la gambi ni stacas, mar\cc
chas e kuras.
\end{VUpage}

% --- Feuillet 15 (folio 13) ------------------------------------------
\begin{VUpage}[15]{13}
%%K t02-15-1 p
\VUblancAlinea
15. --- La \VUgras{shultro} juntas la brakio al korpo.\nl
Tre robusta \VUgras{muskulo,} la bicepso, movas la bra\cc
kio quan la \VUgras{kudo} juntas al \VUgras{avanbrakio.} La\nl
\VUgras{karpo} konstitucas l'artiko di la \VUgras{manuo,} qua\nl
finas per la \VUgras{fingri} e l'\VUgras{ungli.} Sur la manuo ni\nl
dicernas \VUgras{veino} (ed arterio) en qua fluas la\nl
\VUgras{sango.}

%%K t02-16-1 p
\VUblancAlinea
16. --- La parti dil gambo esas : la \VUgras{kruro,}\nl
la \VUgras{genuo} e la \VUgras{poplito,} la \VUgras{suro,} di qua la \VUgras{karno}\nl
esas kovrita per la \VUgras{pelo,} la \VUgras{maleolo} e la \VUgras{pedo.}\nl
La gambal \VUgras{osti} esas tre grosa e tre forta. Ye\nl
l'extremajo dil pedo esas la \VUgras{ped-fingri.} La ce\cc
tera parti esas la \VUgras{plando} e la \VUgras{talono.}

%%K t02-17-1 p
\VUblancAlinea
17. --- Tal esas la preske kompleta deskripto\nl
dil korpo homala.

%%K t02-17-2 p
\VUblancAlinea
\textit{Noto.} --- \textit{Per helpo dil expresuro : « me do}\cc
\textit{loras ye... (la kapo e. c.) » la docisto povos}\nl
\textit{riuzar ica tota vortaro ye la vidpunto di la bona}\nl
\textit{o mala} \VUgras{stando korpala.}

%%K t02-tit-4 sub
\VUsaut{5.0mm}
\VUpk{10.2pt}{40}{Gimnastiko.}
\VUsaut{2.4mm}
\VUcentre{10.2pt}{0}{\textit{(Naraco da un de la lernanti.)}}
\VUsaut{3.0mm}

%%K t02-18-1 p
\VUblancAlinea
18. --- Por esar flexebla, ajila e bonstanda,\nl
ni mustas exercar nia gambi e nia brakii. Pro\nl
to ni ludas e praktikas \VUgras{gimnastiko.}

%%K t02-19-1 p
\VUblancAlinea
19. --- Dum la semano, ica du exerci even\cc
tas en la korto dil liceo (videz la tabelo N\textsuperscript{o} 1).\nl
Ma jovdie e sundie ni iras sur granda gazo\cc
neyo, ube ni havas spaco por kurar ed amuzar\nl
ni.
\end{VUpage}

% --- Feuillet 16 (folio 14) ------------------------------------------
\begin{VUpage}[16]{14}
%%K t02-20-1 p
\VUblancAlinea
20. --- Nia mastri e nia parenti kelkafoye\nl
akompanas ni ibe; ma \VUgras{docisto pri gimnas}\cc
\VUgras{tiko}\textsuperscript{(12)} ya direktas l'exerci.

%%K t02-21-1 p
\VUblancAlinea
21. --- Sur la gazoneyo, sinistre dil enireyo,\nl
trovesas la \VUgras{gimnastikal aparati}\textsuperscript{(1)}, ni klimas\nl
per la \VUgras{skalo}\textsuperscript{(2)}, ni renverse agas ye la \VUgras{ringi}\textsuperscript{(3)}\nl
e ye la \VUgras{trapezo}\textsuperscript{(4)}, ni hisas ni per la manui\nl
til la suprajo di la \VUgras{korda skalo}\textsuperscript{(5)} o di la\nl
\VUgras{kordo nodoza}\textsuperscript{(6)}.

%%K t02-22-1 p
\VUblancAlinea
22. --- Dume, nia kamaradi saltas super la\nl
\VUgras{ligna kavalo}\textsuperscript{(7)} o la \VUgras{stangi paralela}\textsuperscript{(11)}. Altri\nl
\VUgras{boxas}\textsuperscript{(8)} o \VUgras{skermas}\textsuperscript{(9)} kun maskilo e per \VUgras{fle}\cc
\VUgras{reto.} Fine, altri \VUgras{halteragas}\textsuperscript{(10)}.

%%K t02-tit-5 sub
\VUsaut{4.6mm}
\VUpk{10.2pt}{40}{La Ludi.}
\VUsaut{3.0mm}

%%K t02-23-1 p
\VUblancAlinea
23. --- Cirkum la gimnastikanti, pueruli e\nl
puerini ludas per diversa \VUgras{ludi.} La puerini amu\cc
zas su per sua \VUgras{ringego}\textsuperscript{(14)}; eli saltas per \VUgras{kor}\cc
\VUgras{do}\textsuperscript{(13)} od invitas sua fratuli e kuzuli a \VUgras{kro}\cc
\VUgras{ket}\textsuperscript{(15)}-partio. Eli charmesas en forpulsar la\nl
\VUgras{bulo} di l'adverso per sua \VUgras{ligna marteleto,} ye\nl
l'instanto kande lu pensas trairar facile l'ar\cc
keto !

%%K t02-24-1 p
\VUblancAlinea
24. --- La pueruli preferas la ludi plu mus\cc
kulala. Volunte li ludas per \VUgras{kegli}\textsuperscript{(16)} quin li\nl
abatas habile per \VUgras{bulo}\textsuperscript{(17)}. Li anke ne despri\cc
zas ludar per \VUgras{tupio}\textsuperscript{(18)} e \VUgras{bilboketo}\textsuperscript{(19)} o mem\nl
\VUgras{ranoludo}\textsuperscript{(20)}. Ma lia preferata ludi esas ya la\nl
\VUgras{buleti}\textsuperscript{(21)} e la \VUgras{mur-baloneto}\textsuperscript{(22)}, nam la muro\nl
di la \VUgras{teplico}\textsuperscript{(26)} esas tre apta por risaltigar la\nl
lejera baloneto.

%%K t02-25-1 p
\VUblancAlinea
25. --- La mikra Ioannes, iranta per sua\nl
\VUgras{stelti}\textsuperscript{(24)}, esas tre habila e tre bone savas kon\cc
\parplein
\end{VUpage}

% --- Feuillet 17 (folio 15) -------------------------------------------
%  Cette page porte la note de la tabelo 2 : l'appel est un asterisque
%  entre parentheses place apres shultro-salto (30), et la note precise
%  la nomizado en Ido di la infantal e pueral ludi.
\begin{VUpage}[17]{15}
%%K t02-noto-1 noto
\VUnotes{143.2mm}{%
(*) La infantal o pueral ludi havas ofte nomi pure na\cc
cionala. Ni agis por nomizar li en Ido maxim bone posible,\nl
sive inspirante ni per la ago ipsa, per qua li esas karakte\cc
rizata, sive selektante la nacional nomo en ica od ita landi,\nl
kande lu ne esas tro nejusta. Ex. : la \VUgras{barelo-ludo} (Ger\cc
mane e France) esas la \VUgras{rano-ludo} \textsuperscript{(20)} (Angle) en olqua\nl
la ludanti esforcas jetar sua ronda disketi tra la boko di\nl
la metal rano qua beante stacas sur la kesto (barelo) trui\cc
zita.\nl
La \VUgras{shultro-salto}\textsuperscript{(30)} diferas de la \VUgras{dorso-salto} nur per\nl
ico : por saltar super la kapo, la luderi apogas la manui\nl
sur la shultri di sua partenero, vice ke en la « \VUgras{dorso}\cc
\VUgras{salto} » li apogas la manui nur sur ilua dorso. --- Od anke\nl
kelki de la partoprenanti esas inklinita dum ke l'altri saltos\nl
super lia dorso apoganta la manui sur la shultro, l'unesmi\nl
devas tolerar la shoko sen subflexesar, la duesmi devas\nl
saltar sen perdar l'equilibro (Videz la tabelo ipsa).%
}

%%K t02-25-1 p suite
\VUcontinue
servar l'equilibro. Proxim lu, kelka kamaradi\nl
ludas ye \VUgras{celo-trovo}\textsuperscript{(25)} en ta teplico e dop la\nl
\VUgras{hego.} Altri preferas la \VUgras{frap-divino}\textsuperscript{(28)} o la \VUgras{flug}\cc
\VUgras{balono}\textsuperscript{(29)} qua flugetas de un \VUgras{raketo} ad l'altra.

%%K t02-26-1 p
\VUblancAlinea
26. --- Meze dil gazoneyo esas du arbori. An\nl
la pedo di un del du, sep od ok pueruli orga\cc
nizis partio di \VUgras{shultro-salto}\textsuperscript{(30)} (*). Ne en\nl
omna landi konocesas (esas konocata) ta ludo;\nl
ma omni savas quon esas la \VUgras{dorso-salto,} per\nl
qua la pueri ne ludas tafoye.

%%K t02-tit-6 sub
\VUsaut{5.0mm}
\VUpk{10.2pt}{40}{La Ludi en liber aero.}
\VUsaut{3.4mm}

%%K t02-27-1 p
\VUblancAlinea
27. La \VUgras{blindo-ludo}\textsuperscript{(31)} esas anke tre amu\cc
zanta; puerini e pueruli alterne metas la \VUgras{okul}\cc
\VUgras{bendo} an sua okuli e sizas la nehabili qui lasas\nl
su kaptesar.

%%K t02-28-1 p
\VUblancAlinea
28. --- Meze dil gazoneyo, yen ludo tre\nl
Franca ma kelke violentoza : to esas la \VUgras{ursi}\cc
\VUgras{no}\textsuperscript{(32)} multe plu bruisanta kam la tante tran\cc
\parplein
\end{VUpage}

% --- Feuillet 18 (folio 16) --------------------------------------------
\begin{VUpage}[18]{16}
%%K t02-28-1 p suite
\VUcontinue
quila ludo dil \VUgras{kaito}\textsuperscript{(33)} o mem kam la Angla\nl
ludo \VUgras{teniso}\textsuperscript{(34)}.

%%K t02-29-1 p
\VUblancAlinea
29. --- Ta lasta sporto esas la joyo di la\nl
granda lernanti. Nia docisti ipsa volunte prak\cc
tikas olu. Ol postulas granda habileso kun aji\cc
leso ed ol tre plezas a me. Me lasas a la pue\cc
rini la ludi di \VUgras{ocililo}\textsuperscript{(35)}.

%%K t02-30-1 p
\VUblancAlinea
30. --- Me ne prizas altra ludo Angla qua for\cc
san divenus danjeroza.

%%K t02-30-1b p
\VUblancAlinea
Me intencas parolar pri la \VUgras{ped-balono}\textsuperscript{(36)}\nl
qua joyigas mea seniora fratulo. Me preferas\nl
nia olda \VUgras{baro-ludo}\textsuperscript{(38)} tam bone higienala e\nl
vere min brutala.

%%K t02-tit-7 sub
\VUsaut{4.6mm}
\VUpk{10.2pt}{40}{Biciklago e balonludo.}
\VUsaut{3.0mm}

%%K t02-31-1 p
\VUblancAlinea
31. --- Anke volunte me cirkumiras la gazo\cc
neyo per \VUgras{biciklo}\textsuperscript{(37)}.

%%K t02-32-1 p
\VUblancAlinea
32. --- Dum la proxima yaro me lernos la\nl
\VUgras{kavalko}\textsuperscript{(39)}. Tante belesas kavalo qua gracioze\nl
pazas o trotas, e qua galope saltas super\nl
l'obstakli !

%%K t02-33-1 p
\VUblancAlinea
33. --- Somere, ni lernas la \VUgras{natarto}\textsuperscript{(40)}. Ni\nl
delektante plunjas e balnas en la klar e fresh\nl
aquo di la rivero. Kelkafoye fine ni lansas\nl
papera \VUgras{balono}\textsuperscript{(41)} qua majestoze acensas en\nl
l'atmosfero, quik kande ol esas plenigita per\nl
varma aero.

%%K t02-34-1 p
\VUblancAlinea
34. --- Esas anke mult altra ludi, ma me ne\nl
finus enumerar oli ed on koncias segun ica\nl
rezumo, ke en la jovdii, on ne tre enoyas sur\nl
nia bela gazoneyo.

%%K t02-34-2 p
\VUblancAlinea
N.-B. --- \textit{La docisto facile kompletigos ta cha}\cc
\textit{pitro; plu detaloze deskriptante singla de la}\nl
\textit{maxim importanta ludi.}
\end{VUpage}
